% GENERATED FILE. Edit canonical lessons, then run npm run generate:books.
% canonical-source-hash: fnv1a64:5ba8a3fbdabab32f
% canonical-lessons: GU-C319-killo, GU-C319-khander, GU-C319-smarak, GU-C319-murti, GU-C319-vasvu

\chapter{The Old Fort}
\label{ch:gu-the-old-fort}
\hlchaptermodality{319}

\begin{chapteropening}
\textbf{By the end of this chapter:} \emph{I can describe an old fort and what remains of it.}

Complete the last lesson of chapter 319: I can describe an old fort and what remains of it.
\end{chapteropening}

\section[\texorpdfstring{\gu{કિલ્લો} (killo)}{કિલ્લો (killo)}]{\gu{કિલ્લો} (killo) — a fort}

\label{lesson:GU-C319-killo}



\begin{quote}
Before the new one: say the Gujarati for a royal court, then the Gujarati for a throne.
\end{quote}

\subsection*{You'll want to know: \gu{કિલ્લો}}
\textbf{\gu{કિલ્લો}} — \emph{killo} — \textquotedblleft{}a fort\textquotedblright{}.

\begin{quote}
\gu{ટેકરી} \gu{પર} \gu{એક} \gu{જૂનો} \gu{કિલ્લો} \gu{છે}.
\end{quote}

\begin{quote}
\emph{ṭekrī par ek jūno killo che.}
\end{quote}

\begin{quote}
There is an old fort on the hill.
\end{quote}

A masculine noun.

\begin{cousinweb}
\textbf{\gu{કિલ્લો}} (\emph{killo}) is usually traced, through Persian, to Arabic \emph{qalʿa}, \textquotedblleft{}a fortress\textquotedblright{}. It is masculine: \textbf{\gu{જૂનો} \gu{કિલ્લો}}, an old fort.
\end{cousinweb}

\begin{grammarlens}[title={What you've built}]
One more word for telling what happened.
\end{grammarlens}

\subsection*{Your turn}
\begin{itemize}
\raggedright
  \item \emph{Say it:} \emph{killo}
  \item \emph{Say it:} \emph{killo}, once more
  \item \emph{Say it:} \emph{ṭekrī par ek jūno killo che.} — There is an old fort on the hill.
  \item \emph{Recall:} say the Gujarati for a royal court, then the Gujarati for a throne, then say \emph{killo} again
\end{itemize}

\begin{tcolorbox}[breakable,colback=teal!4,colframe=teal!35!black,title={Before you move on}]
What does \gu{કિલ્લો} mean? (\textquotedblleft{}A fort\textquotedblright{}.) Say it once more.
\end{tcolorbox}
\section[\texorpdfstring{\gu{ખંડેર} (khanḍer)}{ખંડેર (khanḍer)}]{\gu{ખંડેર} (khanḍer) — ruins}

\label{lesson:GU-C319-khander}



\begin{quote}
Before the new one: say the Gujarati for a throne, then the Gujarati for a fort.
\end{quote}

\subsection*{You'll want to know: \gu{ખંડેર}}
\textbf{\gu{ખંડેર}} — \emph{khanḍer} — \textquotedblleft{}ruins\textquotedblright{}.

\begin{quote}
\gu{આજે} \gu{ત્યાં} \gu{ફક્ત} \gu{ખંડેર} \gu{છે}.
\end{quote}

\begin{quote}
\emph{āje tyā̃ phakt khanḍer che.}
\end{quote}

\begin{quote}
Today only ruins remain there.
\end{quote}

A neuter noun.

\begin{cousinweb}
\textbf{\gu{ખંડેર}} (\emph{khanḍer}) probably goes back to \emph{*khaṇḍitaghara}, \textquotedblleft{}a broken house\textquotedblright{}, a form CDIAL reconstructs, where it lists Gujarati \emph{khãḍiyer}, \textquotedblleft{}ruins\textquotedblright{} (CDIAL 3798). It is what is left when walls and roofs have fallen.
\end{cousinweb}

\begin{grammarlens}[title={What you've built}]
One more word for telling what happened.
\end{grammarlens}

\subsection*{Your turn}
\begin{itemize}
\raggedright
  \item \emph{Say it:} \emph{khanḍer}
  \item \emph{Say it:} \emph{khanḍer}, once more
  \item \emph{Say it:} \emph{āje tyā̃ phakt khanḍer che.} — Today only ruins remain there.
  \item \emph{Recall:} say the Gujarati for a throne, then the Gujarati for a fort, then say \emph{khanḍer} again
\end{itemize}

\begin{tcolorbox}[breakable,colback=teal!4,colframe=teal!35!black,title={Before you move on}]
What does \gu{ખંડેર} mean? (\textquotedblleft{}Ruins\textquotedblright{}.) Say it once more.
\end{tcolorbox}
\section[\texorpdfstring{\gu{સ્મારક} (smārak)}{સ્મારક (smārak)}]{\gu{સ્મારક} (smārak) — a memorial}

\label{lesson:GU-C319-smarak}



\begin{quote}
Before the new one: say the Gujarati for a fort, then the Gujarati for ruins.
\end{quote}

\subsection*{You'll want to know: \gu{સ્મારક}}
\textbf{\gu{સ્મારક}} — \emph{smārak} — \textquotedblleft{}a memorial\textquotedblright{}.

\begin{quote}
\gu{ગામના} \gu{લોકોએ} \gu{ત્યાં} \gu{એક} \gu{સ્મારક} \gu{બનાવ્યું}.
\end{quote}

\begin{quote}
\emph{gāmnā lokoe tyā̃ ek smārak banāvyũ.}
\end{quote}

\begin{quote}
The villagers built a memorial there.
\end{quote}

A neuter noun.

\begin{cousinweb}
\textbf{\gu{સ્મારક}} (\emph{smārak}) is a tatsama, Sanskrit \emph{smāraka}, \textquotedblleft{}recalling, reminding of\textquotedblright{} (Monier-Williams): a thing put up so that people remember.
\end{cousinweb}

\begin{grammarlens}[title={What you've built}]
One more word for telling what happened.
\end{grammarlens}

\subsection*{Your turn}
\begin{itemize}
\raggedright
  \item \emph{Say it:} \emph{smārak}
  \item \emph{Say it:} \emph{smārak}, once more
  \item \emph{Say it:} \emph{gāmnā lokoe tyā̃ ek smārak banāvyũ.} — The villagers built a memorial there.
  \item \emph{Recall:} say the Gujarati for a fort, then the Gujarati for ruins, then say \emph{smārak} again
\end{itemize}

\begin{tcolorbox}[breakable,colback=teal!4,colframe=teal!35!black,title={Before you move on}]
What does \gu{સ્મારક} mean? (\textquotedblleft{}A memorial\textquotedblright{}.) Say it once more.
\end{tcolorbox}
\section[\texorpdfstring{\gu{મૂર્તિ} (mūrti)}{મૂર્તિ (mūrti)}]{\gu{મૂર્તિ} (mūrti) — a statue; an idol}

\label{lesson:GU-C319-murti}



\begin{quote}
Before the new one: say the Gujarati for ruins, then the Gujarati for a memorial.
\end{quote}

\subsection*{You'll want to know: \gu{મૂર્તિ}}
\textbf{\gu{મૂર્તિ}} — \emph{mūrti} — \textquotedblleft{}a statue; an idol\textquotedblright{}.

\begin{quote}
\gu{કિલ્લામાં} \gu{પથ્થરની} \gu{એક} \gu{સુંદર} \gu{મૂર્તિ} \gu{છે}.
\end{quote}

\begin{quote}
\emph{killāmā̃ paththarnī ek sundar mūrti che.}
\end{quote}

\begin{quote}
There is a beautiful stone statue in the fort.
\end{quote}

A feminine noun.

\begin{cousinweb}
\textbf{\gu{મૂર્તિ}} (\emph{mūrti}) is a tatsama, Sanskrit \emph{mūrti}, \textquotedblleft{}a solid body, a form, an image\textquotedblright{} (Monier-Williams). It is a statue of a person, and in a temple the image of a god.
\end{cousinweb}

\begin{grammarlens}[title={What you've built}]
One more word for telling what happened.
\end{grammarlens}

\subsection*{Your turn}
\begin{itemize}
\raggedright
  \item \emph{Say it:} \emph{mūrti}
  \item \emph{Say it:} \emph{mūrti}, once more
  \item \emph{Say it:} \emph{killāmā̃ paththarnī ek sundar mūrti che.} — There is a beautiful stone statue in the fort.
  \item \emph{Recall:} say the Gujarati for ruins, then the Gujarati for a memorial, then say \emph{mūrti} again
\end{itemize}

\begin{tcolorbox}[breakable,colback=teal!4,colframe=teal!35!black,title={Before you move on}]
What does \gu{મૂર્તિ} mean? (\textquotedblleft{}A statue; an idol\textquotedblright{}.) Say it once more.
\end{tcolorbox}
\section[\texorpdfstring{\gu{વસવું} (vasvũ)}{વસવું (vasvũ)}]{\gu{વસવું} (vasvũ) — to settle (in a place)}

\label{lesson:GU-C319-vasvu}



\begin{quote}
Before the new one: say the Gujarati for a memorial, then the Gujarati for a statue.
\end{quote}

\subsection*{You'll want to know: \gu{વસવું}}
\textbf{\gu{વસવું}} — \emph{vasvũ} — \textquotedblleft{}to settle (in a place)\textquotedblright{}.

\begin{quote}
\gu{ઘણાં} \gu{વરસ} \gu{પહેલાં} \gu{લોકો} \gu{આ} \gu{નદી} \gu{કિનારે} \gu{વસ્યા}.
\end{quote}

\begin{quote}
\emph{ghaṇā̃ varas pahelā̃ loko ā nadī kināre vasyā.}
\end{quote}

\begin{quote}
Many years ago people settled on the bank of this river.
\end{quote}

Past (he …): \gu{તે} \gu{અહીં} \gu{વસ્યો}, \emph{te ahī̃ vasyo}.

\begin{cousinweb}
\textbf{\gu{વસવું}} (\emph{vasvũ}) is inherited from Sanskrit \emph{vásati}, \textquotedblleft{}dwells\textquotedblright{}, and CDIAL lists the Gujarati verb, \textquotedblleft{}settles\textquotedblright{} (CDIAL 11435). People \textbf{\gu{વસે} \gu{છે}} where they make their homes; \textbf{\gu{રહેવું}} is everyday living.
\end{cousinweb}

\begin{grammarlens}[title={What you've built}]
One more word for telling what happened.
\end{grammarlens}

\subsection*{Your turn}
\begin{itemize}
\raggedright
  \item \emph{Say it:} \emph{vasvũ}
  \item \emph{Say it:} \emph{vasvũ}, once more
  \item \emph{Say it:} \emph{ghaṇā̃ varas pahelā̃ loko ā nadī kināre vasyā.} — Many years ago people settled on the bank of this river.
  \item \emph{Recall:} say the Gujarati for a memorial, then the Gujarati for a statue, then say \emph{vasvũ} again
\end{itemize}

\begin{tcolorbox}[breakable,colback=teal!4,colframe=teal!35!black,title={Before you move on}]
What does \gu{વસવું} mean? (\textquotedblleft{}To settle (in a place)\textquotedblright{}.) Say it once more.
\end{tcolorbox}
